\documentclass[11pt]{article}
\usepackage{amsmath, amssymb, amsthm}
\usepackage[margin=1in]{geometry}
\usepackage[pdftitle={Day 170 - Theorem B PROVED. Year-arc terminates.},pdfauthor={grandpa-rick}]{hyperref}
\usepackage{parskip}
\usepackage{booktabs}

\newcommand{\R}{\mathcal{R}}
\newcommand{\barD}{\bar D}
\newcommand{\SOURCE}{\mathrm{SOURCE}}

\title{Day 170 --- Theorem B PROVED. Year-arc terminates. \\
\small $\barD|_{E_3=0}$ closed form proved unconditionally as a ring identity in $\mathbb{Q}(T,s,p)[Y]/(pTY^2+(sT-1)Y+T)$}
\author{Rick (grandpa-rick, work-in-progress)}
\date{2026-09-05}

\begin{document}
\maketitle

\noindent\textbf{Header stamp.}\\
Author: Rick (grandpa-rick). \quad Date: 2026-09-05 (Day 170). \\
Recipient: Clio Vega. \quad CC: Robin Langer. \\
Title: Theorem B PROVED. Year-arc terminates. \\
Source commit (work-in-progress): \texttt{db21340}.

\bigskip

\noindent\textbf{Status.} MAJOR WIN. The three-way collapse class from Day 165 ($\Sigma_0 \Longleftrightarrow R^{(-1)} \Longleftrightarrow$ Theorem B) is now \texttt{proved}. Missing Lemma (R) is closed unconditionally. C.5 (the $E_3=0$ layer identity from Day 156, equivalently Day 161 Theorem 4) auto-promotes to \texttt{proved}. \textbf{FPSAC \S 5 open list of structural theorems: $1 \to 0$.}

The year-long $b_k / \Psi / P_b / \mathrm{C.5}$ arc that started with $b_k \equiv 0 \pmod 3$ around Day 120 and threaded through Days 148 (Horn hypergeometric), 152 ($\psi$ closed form), 154 (Narayana at $E_3=0$), 156 (C.5 conjectured), 158--169 (Riccati layers) \textbf{terminates here.}

\section{The result}

\noindent\textbf{Theorem B (Day 162, now PROVED).} In the algebraic ring
$$\R \;:=\; \mathbb{Q}(T, s, p)[Y]/(pTY^2 + (sT-1)Y + T), \qquad q \;=\; 1 - sT - 2pTY,$$
one has
$$\boxed{\;\barD\big|_{E_3=0} \;=\; \frac{T Y^2\bigl[(q+1)^2 - E_1 T\bigr]}{q^3}.\;}$$

Two equivalent formulations (both promoted from \texttt{checked-sober} to \texttt{proved}):

\begin{itemize}
  \item \textbf{$R^{(-1)}$ closed form (Day 162):}
  $$R^{(-1)} \;=\; \frac{T\bigl[E_2 Y^2\bigl((q+1)^2 - E_1 T\bigr) + (q + R_1 R_2)/2\bigr]}{q^3}, \qquad R_1 R_2 = 1 - T^2(E_1^2 - 4 E_2).$$
  \item \textbf{$\Sigma_0$ closed form (Day 165):}
  $$-\Sigma_0 \;=\; \frac{(q + 1 - u)(q^2 - 6q + 6 - 6u)}{2\, q^4}, \qquad u = E_1 T.$$
\end{itemize}

All three are pairwise equivalent by Day 165's three-way collapse; proving one proves all.

\section{The proof (30-second version)}

Combine four earlier results:

\begin{enumerate}
  \item \textbf{Prop 3 (Day 167, PROVED).} For all $n$,
  $$R^{(-1)}_n \;=\; \tfrac{1}{2}\,\partial_{u_3}^2 \Xi_n\big|_{u_3=0} \;-\; [\deg_{(u_1,u_2)} = n{-}1]\bigl([T^n]\log(F_{-1}/F_0)\bigr),$$
  proved by weight-grading (no constructive $F_{-1}$ needed).

  \item \textbf{Route A closed form (Day 167, PROVED).} $\tfrac{1}{2}\partial_{u_3}^2 \Xi|_{u_3=0}$ has an explicit closed form in $\xi_0, \xi_1, \xi_2$ via the chain rule at $u_3 = 0$.

  \item \textbf{$L_0$ closed form (Day 168, PROVED).} The sub-sub-top layer of $G_0 := F_0'/F_0$ equals
  $$L_0 \;=\; \frac{1 + 3 T K_0 + T^2 K_0^2 + T \theta K_0}{q}, \qquad K_0 \;=\; \frac{p Y (2q+1) + s q}{q^2}.$$

  \item \textbf{$L_{-1}$ closed form (Day 169, PROVED, corrected).} The sub-sub-top layer of $G_{-1} := F_{-1}'/F_{-1}$ equals $L_{-1} = -\SOURCE/(q^3 H)$ with the corrected $\SOURCE$ (see \S 3). Equivalently, in $\R$,
  $$L_{-1} \;=\; \frac{A_0 + A_1\, q + (B_0 + B_1\, q)\, Y}{Y\, q^5}.$$
\end{enumerate}

Set $I_\Delta := \int_0^T (L_{-1} - L_0)\, dT'$. By Prop 3,
\begin{equation}\tag{P3}
R^{(-1)}(T) \;=\; \tfrac{1}{2}\partial_{u_3}^2 \Xi\big|_0 \;-\; I_\Delta(T).
\end{equation}
Both sides vanish at $T=0$; their $T$-derivatives satisfy
\begin{equation}\tag{P3$'$}
\partial_T R^{(-1)} \;+\; (L_{-1} - L_0) \;=\; \partial_T\Bigl[\tfrac{1}{2}\partial_{u_3}^2 \Xi\big|_0\Bigr].
\end{equation}
Under $q = 1 - sT - 2pTY$, the RHS minus LHS is a rational function in $\{T, s, p, Y\}$. Reducing modulo $pTY^2 + (sT-1)Y + T$ via \texttt{sp.div} + \texttt{sp.subs} + \texttt{sp.cancel} gives $0$ in $\R$ (runtime $\approx 0.5$ s). Both sides vanish at $T=0$, so equality follows. $\square$

\section{Correction to Day 169's SOURCE formula}

The Day 169 proof file (\S 3.3) lists the $\SOURCE$ for $L_{-1}$ as a sum of 12 terms. \textbf{One term is missing}: the $P_3 G^3$ contribution at layer $e=1$ (the $H^2 K$ part of the $\delta=2$ diagonal), which contributes $18\, T^3\, H^2\, K$.

The correct $\SOURCE$ is:
\begin{align*}
\SOURCE(T) \;=\;& R_3 H'' + [-11T + 14sT^2 + (12p - 3s^2) T^3]\, H' \\
&+ [1 + 12 s T + (5 p - s^2) T^2]\, H \\
&+ 3 R_3\, (H K' + K H') \\
&+ [23 T^2 + s T^3]\, H^2 + 18 T^3\, H H' + T^4\, H^3 + 3 R_3\, H K^2 \\
&+ R_2 K' + 2\,[-11T + 14 s T^2 + (12p - 3s^2) T^3]\, H K \\
&+ [-s + (2s^2 + 10 p) T + (4 p s - s^3) T^2]\, K + R_2 K^2 \\
&+ \boxed{\;18\, T^3\, H^2\, K\;} \qquad \text{(missing from Day 169's writeup)},
\end{align*}
with $H = pY/T$, $K = K_{-1} = -pY/q^2$, $R_3 = -T^2 q^2$, $R_2 = q^2(1 - sT)$.

\emph{Trace of the error.} Day 169's step 15 (\texttt{step15\_L\_closed\_form.py}) correctly enumerates the $P_3 G^3$ contributions at $e=0,1,2$; step 16 (\texttt{step16\_solve\_L.py}) correctly includes \texttt{c\_18T3\_H2K} in the assembly. The proof-file writeup omitted the term in transcription. Numerics using step 16 gave correct $L_{-1}$ series, so the numerical check was clean --- only the human writeup was incomplete. The Day 170 discipline of re-running the checker script against the writeup caught the omission. (See feedback rule: \emph{never trust the writeup, only the running code}.)

\emph{Compact ring form.} Using $q = 1 - sT - 2pTY$ and $Y^2 = ((1-sT)Y - T)/(pT)$, along with $q^2 = (1-sT)^2 - 4 p T^2$, the corrected $L_{-1}$ reduces in $\R$ to
$$L_{-1} \;=\; \frac{A_0 + A_1\, q + (B_0 + B_1\, q)\, Y}{Y\, q^5}$$
with
\begin{align*}
A_0 &= -T(4 T^2 p - T^2 s^2 + 4 T s - 3), & A_1 &= 24 T,\\
B_0 &= -4 T^3 p s + T^3 s^3 + 41 T^2 p - 15 T^2 s^2 + 27 T s - 13, & B_1 &= 14(T s - 1).
\end{align*}
Note $(4 T^2 p - T^2 s^2 + 2 T s - 1)^2 = q^4$, so the denominator is $Y q^5$ as claimed.

\section{Machine verification}

Location: \texttt{scripts/day170/} in this repository.

\begin{itemize}
  \item \texttt{step13\_Lm1\_corrected\_SOURCE.py} --- verify $L_{-1}$ series with 13-term $\SOURCE$ matches \texttt{FP\_coeffs} to $n=8$. PASS.
  \item \texttt{step17\_full\_proof\_check.py} --- independent numerical cross-checks at $(s, p) \in \{(2,3), (1,1), (5,2)\}$: all $T$-coefficients of $\mathcal{R}\!A$ and $\mathcal{L}$ agree for $n \le 9$. 10/10 at each.
  \item \texttt{step18\_clean\_proof.py} --- clean algebraic proof. Prints \texttt{"num reduced (poly in Y, deg <= 1): 0"} after $\approx 0.5$ s of \texttt{sp.cancel} + \texttt{sp.subs} + \texttt{sp.div}. This is the file that discharges the theorem.
\end{itemize}

The proof file \texttt{proofs/2026-09-05-day170-theorem-B-PROVED.md} \S 4 contains the full derivation; the collaborator draft in \texttt{memory/for-collaborator/2026-09-05-day170-theorem-B-proved.md} is the source of the present note.

\section{Boundary at $T=0$}

\begin{itemize}
  \item $\tfrac{1}{2}\partial_{u_3}^2 \Xi|_{u_3=0}$ is $O(T)$ (since $\log F_P = O(T)$).
  \item $R^{(-1)}$ (Day 162 closed form) has an explicit factor of $T$ in the numerator.
  \item $I_\Delta = \int_0^T (L_{-1} - L_0)\, dT'$ vanishes at $T=0$.
\end{itemize}
Hence $\bigl[\tfrac{1}{2}\partial_{u_3}^2 \Xi|_0\bigr] - I_\Delta - R^{(-1)}$ vanishes at $T = 0$ and has zero $T$-derivative in $\R$, so equals $0$ as a formal power series.

\section{Registry updates}

\begin{center}
\begin{tabular}{@{}lll@{}}
\toprule
Node & Before Day 170 & After Day 170 \\
\midrule
\texttt{bar-D-closed-form-E3-zero} (Theorem B) & \texttt{checked-sober} ($n \le 14$) & \textbf{\texttt{proved}} \\
\texttt{R-minus-one-closed-form} (Day 162) & \texttt{checked-sober} ($n \le 14$) & \textbf{\texttt{proved}} \\
\texttt{LA-F1-sub-top-Sigma-0} (Day 165) & \texttt{checked-sober} ($n \le 24$, 15 specs) & \textbf{\texttt{proved}} \\
\texttt{narayana-layer-d1-E3-zero} (C.5) & \texttt{computed} & \textbf{\texttt{proved}} \\
\texttt{Missing-Lemma-R} (implicit) & \texttt{proved} conditional on Thm B & \textbf{\texttt{proved} unconditionally} \\
\texttt{day170-prop3-ring-identity} (NEW) & --- & \textbf{\texttt{proved}} \\
\bottomrule
\end{tabular}
\end{center}

\section{What Clio might want to know}

Route (v) --- your reduction of Missing Lemma (R) via Prop 3 --- is now proved end-to-end. The chain $\{$Prop 3, Route A, Route B$\}$ that you outlined is the exact structure of the Day 170 proof: Prop 3 is Day 167 (weight-grading), Route A is Day 167, Route B is Days 168+169 (Riccati of the $F_0$ and $F_{-1}$ ODEs). Day 170 itself is only the ring identity that glues them.

Your $E_2$-shift conjecture (verified 26/26 for $(n, b) \in \{4, \ldots, 7\} \times \{0, \ldots, 6\}$ on Day 169) stays at \texttt{computed}. The shift-law proof is a separate port of Day 131 \S 4 with $3 \to n$ and is on the Day 172+ queue. It is not part of this closure.

\section{What Robin might want to know}

The FPSAC 2027 abstract can now be drafted with a full theorem statement, complete proof, and no \texttt{checked-sober} on the crown result. Deadline is estimated late March / early April 2027 (Browse 129: the FPSAC 2027 dates page 404'd; extrapolating from 2026). Chairs are D'Adderio, Pilaud, Rajchgot. Invited speakers include Haiman and Martha Yip; Mishna (kernel-method / BM\&J expert) is on the invitee list per Browse 128 --- the natural audience for the algebraic-GF approach.

Community context (Browse 129): Griffin--Mellit \texttt{2504.06936} (14 citations in 5 months) is the new dominant hub; Hikita's SS proof (\texttt{2410.12758}) is the anchor. Nobody has an algebraic GF for $X_{P_n}(x; q)$ in the community; Theorem B is first. The natural next question is $q$-polynomial positivity of $e$-coefficients, which nobody has an approach to at the algebraic level.

\section{Post-arc queue (top three)}

\begin{itemize}
  \item Draft FPSAC 2027 abstract (Day 171 PROVE-session output).
  \item Read Tom--Vailaya \texttt{2503.19344}: does gluing already give $e$-positivity for $P_n$ as a corollary? Adjust framing accordingly.
  \item 10-line sympy check: does the closed form for $\psi$ at $t=0$ reproduce Kim--Lee--Yoo's Hall--Littlewood expansion? Would subsume their result as a corollary.
\end{itemize}

\end{document}
